\documentclass[11pt,a4paper]{coverletter}

\senderblock
  {Radheem Bin Razi}
  {Distributed Systems \& Cloud-Native Engineer}
  {Ilmenau, Germany \textbar\ \href{mailto:sheikh.radheem@gmail.com}{sheikh.radheem@gmail.com} \textbar\ \href{https://radheem.github.io/my-cv}{portfolio}}

\begin{document}

%% ===== ENGLISH =====
\recipient{Fraunhofer}{Hiring Team}{}

\opening{Dear Hiring Team,}

Modern maritime logistics depends on turning scattered sensor feeds into a single source of truth before a vessel even docks. Building that quiet kind of reliability across heterogeneous data streams is the part of data engineering I enjoy most. I am applying for the Wissenschaftliche Hilfskraft position in Data Lakehouse \& Data Engineering at the Fraunhofer CML.

My recent work has focused on stitching together real-time telemetry and batch archives into consistent, queryable stores. I designed a Kafka-driven metrics pipeline that ingested per-device measurements and fanned them out to time-series and document databases. The system handled continuous streams without dropping records. I configured the consumer logic to align with strict data-quality gates. This experience maps directly to your goal of connecting live AIS, radar, and camera feeds to a local Lakehouse. I am comfortable configuring open-source streaming and storage layers. I prefer building pipelines that survive network hiccups and worker restarts.

I have also orchestrated complex ingestion workflows that require reliable scheduling and transparent state tracking. I built a five-step document-processing pipeline that extracts, chunks, embeds, summarizes, and indexes content while streaming progress updates. Each step included configurable timeouts and automatic retries. The overall job kept running even when individual workers failed. I paired these pipelines with declarative routing and automated validation checks. These checks caught schema drift before it reached production. This approach aligns with your need for robust ETL processes and self-directed literature research to evaluate emerging Lakehouse architectures. I thrive in research environments where technical choices must balance experimental flexibility with production-grade stability.

My background in distributed data pipelines and open-source infrastructure will support the CML’s research projects on maritime sensor integration. I am available to start immediately. I can work full-time or in a working-student arrangement. I am currently based in Germany. I can relocate to Hamburg within two to three weeks. A Blue Card is straightforward to process. The employer only needs to issue the employment contract. I look forward to discussing how my pipeline experience can accelerate your maritime data initiatives. Thank you for considering my application.

\closing{Sincerely,}{Radheem Bin Razi}

\clearpage
\selectlanguage{ngerman}

%% ===== DEUTSCH =====
\recipient{Fraunhofer}{Personalabteilung}{}

\opening{Sehr geehrtes Hiring-Team,}

Moderne maritime Logistik hängt davon ab, verteilte Sensor-Streams in eine einzige verlässliche Datenquelle zu verwandeln, noch bevor ein Schiff anlegt. Diese unauffällige, aber verlässliche Stabilität über heterogene Datenströme hinweg aufzubauen, ist der Bereich der Data Engineering, der mir am meisten liegt. Hiermit bewerbe ich mich auf die Position als Wissenschaftliche Hilfskraft im Bereich Data Lakehouse \& Data Engineering am Fraunhofer CML.

Meine jüngsten Projekte konzentrierten sich auf das nahtlose Zusammenführen von Echtzeit-Telemetrie- und Batch-Archiven zu konsistenten, abfragefähigen Speichern. Ich entwarf eine Kafka-basierte Metriken-Pipeline, die gerätespezifische Messwerte aufnahm und an Zeitreihen- und Dokumentendatenbanken weiterleitete. Das System verarbeitete kontinuierliche Streams ohne Datenverlust. Ich konfigurierte die Consumer-Logik in Übereinstimmung mit strengen Datenqualitätsprüfungen. Diese Erfahrung passt direkt zu Ihrem Ziel, Live-AIS-, Radar- und Kameradatenströme mit einem lokalen Lakehouse zu verbinden. Ich bin erfahren im Konfigurieren von Open-Source-Streaming- und Speicherschichten. Ich bevorzuge den Aufbau von Pipelines, die auch bei Netzwerkunterbrechungen und Worker-Neustarts stabil laufen.

Zudem habe ich komplexe Ingestion-Workflows orchestriert, die zuverlässiges Scheduling und transparente Statusverfolgung erfordern. Ich entwickelte eine fünfstufige Dokumentenverarbeitungs-Pipeline, die Inhalte extrahiert, chunkt, einbettet, zusammenfasst und indiziert, während Fortschrittsupdates gestreamt werden. Jeder Schritt verfügte über konfigurierbare Timeouts und automatische Wiederholungsversuche. Der Gesamtjob lief weiter, auch wenn einzelne Worker ausfielen. Ich kombinierte diese Pipelines mit deklarativem Routing und automatisierten Validierungschecks. Diese Checks fingen Schema-Drift auf, bevor sie die Produktionsumgebung erreichte. Dieser Ansatz entspricht Ihrem Bedarf an robusten ETL-Prozessen sowie der selbstständigen Literaturrecherche zur Bewertung neu entstehender Lakehouse-Architekturen. Ich entfalte mich besonders gut in Forschungsumgebungen, in denen technische Entscheidungen experimentelle Flexibilität mit produktionsreifer Stabilität in Einklang bringen müssen.

Mein Hintergrund in verteilten Datenpipelines und Open-Source-Infrastrukturen wird die Forschungsprojekte des CML zur maritimen Sensorintegration unterstützen. Ich stehe ab sofort zur Verfügung. Ich kann in Vollzeit oder als Werkstudent arbeiten. Ich befinde mich derzeit in Deutschland. Ein Umzug nach Hamburg ist innerhalb von zwei bis drei Wochen möglich. Die Bearbeitung einer Blue Card ist unkompliziert. Der Arbeitgeber muss lediglich den Arbeitsvertrag ausstellen. Ich freue mich darauf, in einem persönlichen Gespräch zu besprechen, wie meine Erfahrung mit Pipelines Ihre maritimen Dateninitiativen vorantreiben kann. Vielen Dank für die Prüfung meiner Bewerbung.

\closing{Mit freundlichen Grüßen,}{Radheem Bin Razi}

\end{document}
